% !TITLE 春望
% !AUTHOR 杜甫
% !DYNASTY 唐
% !ORDER 21

\begin{poem}
国破山河在，城春草木深。
感时花溅泪，恨别鸟惊心。
烽火连三月，家书抵万金。
白头搔更短，浑欲不胜簪。
\end{poem}

\begin{appreciation}
这是杜甫在安史之乱中被困长安时写的诗。至德二载（757年）的春天，长安已经被叛军占领半年多了，杜甫因为官职太小，来不及逃跑，被留在了这座危城中。

国破山河在，城春草木深——国家破碎了，但山河还在；春天来了，城里的草木却长得很深。在是第一句的关键：国破了，山河还在——这是一种残酷的安慰。城春草木深——长安本是繁华的帝都，如今却荒草丛生。深字写得好，不是浅草，是深草，说明已经很久没有人烟了。

感时花溅泪，恨别鸟惊心——因为感伤时事，看到花开也会流泪；因为怨恨离别，听到鸟叫也会心惊。这两句的妙处在于它们有两种读法：一是花溅泪、鸟惊心，即花像人一样流泪，鸟像人一样心惊；二是感时溅泪、恨别惊心，即诗人看到花而溅泪，听到鸟而惊心。无论哪种读法，都写出了诗人内心的极度敏感。在一个正常人眼中，花开鸟鸣是美好的；但在一个身处乱世的人眼中，一切美好都成了刺痛的提醒。

烽火连三月，家书抵万金——战争已经持续了三个月，一封家书价值万金。抵万金不是夸张，是真实的感受。当生死未卜的时候，知道家人还活着，就是最大的幸福。

白头搔更短，浑欲不胜簪——白发越搔越短，几乎插不住簪子了。搔是抓、挠的意思。人在焦虑时会不由自主地挠头，这是一个下意识的动作。浑欲不胜簪——简直要插不住簪子了。一个老人，在战火中煎熬，头发一把一把地掉。这个细节比任何控诉都更有力量，因为它是一个具体的人、具体的身体在承受具体的痛苦。
\end{appreciation}
